\section{Introduction}

Gesture recognition is a fundamental problem in human-computer interaction (HCI), enabling
intuitive and natural control of machines without physical contact. By interpreting the movements
and positions of the human hand, gesture recognition systems find applications in prosthetic limb
control, rehabilitation robotics, sign language interpretation, virtual reality, and industrial
automation \citep{Phinyomark2012}. Among the available sensing modalities, Electromyography (EMG)
has emerged as one of the most widely adopted approaches, as it directly measures the electrical
activity generated by skeletal muscle contractions during voluntary movement, providing a
physiological signal of motor intent.

EMG signals offer rich temporal and spectral information about the neuromuscular activity
underlying hand gestures. However, EMG-based gesture classification presents several challenges:
signals are inherently noisy, highly subject-dependent, sensitive to electrode placement, and may
exhibit non-stationarities. Despite the growing adoption of deep learning architectures in this
domain, classical machine learning pipelines combining handcrafted feature extraction with
supervised classifiers remain highly relevant due to their interpretability, low computational
requirements, and competitive performance on moderate-sized datasets \citep{Atzori2016}.

This project implements and evaluates a complete machine learning classification pipeline for
distinguishing two hand gestures --- Rock and Paper --- from single-channel EMG signals. Binary
classification was selected to establish a clear performance baseline and systematically
investigate the contribution of individual pipeline components. The pipeline covers all stages:
data loading, missing value analysis, multi-domain feature extraction, statistical feature
selection, PCA dimensionality reduction, comparison of three classifiers under three validation
strategies, a feature-domain ablation study, and a feature importance analysis.
